% =====================================================================
\chapter{Ángulos, ortogonalidad y similitud coseno}
\label{tema:2}
% =====================================================================

La desigualdad de Cauchy--Schwarz permite definir el ángulo entre dos vectores en cualquier dimensión. Este tema desarrolla esa geometría: ortogonalidad, proyección sobre una recta y similitud coseno, que es la medida con la que se comparan los embeddings. Se sigue la sección~1.2 de Strang~\cite{strang} y las secciones~3.3, 3.4 y~3.8 de~\cite{mml}; el comportamiento en alta dimensión procede del capítulo~2 de Blum, Hopcroft y Kannan~\cite{bhk}.

\section{El ángulo entre dos vectores}

\begin{definicion}[Ángulo]\label{def:angulo}
Sean $u,v\in\R^{d}$ no nulos. El \emph{ángulo} entre $u$ y $v$ es el único $\theta\in[0,\pi]$ tal que
\[
\cos\theta=\frac{u\cdot v}{\norm{u}\,\norm{v}}.
\]
\end{definicion}

La definición es correcta porque, por el \cref{teo:cs}, el cociente está en $[-1,1]$, y la función coseno es una biyección de $[0,\pi]$ en $[-1,1]$. La proposición siguiente muestra que, en el plano, coincide con el ángulo de la geometría elemental.

\begin{proposicion}\label{prop:ley-coseno}
Sean $u,v\in\R^{2}$ no nulos y sea $\varphi$ el ángulo geométrico que forman las flechas $u$ y $v$. Entonces $u\cdot v=\norm{u}\norm{v}\cos\varphi$.
\end{proposicion}

\begin{proof}
Los vectores $u$, $v$ y $u-v$ forman un triángulo cuyos lados miden $\norm{u}$, $\norm{v}$ y $\norm{u-v}$, y el ángulo opuesto al lado $u-v$ es $\varphi$. El teorema del coseno da $\norm{u-v}^{2}=\norm{u}^{2}+\norm{v}^{2}-2\norm{u}\norm{v}\cos\varphi$. Por otra parte, la linealidad del producto escalar da $\norm{u-v}^{2}=\norm{u}^{2}+\norm{v}^{2}-2\,u\cdot v$. Igualando ambas expresiones se obtiene el enunciado. (Si $u$ y $v$ son paralelos, el triángulo degenera y la igualdad se comprueba directamente con $u=cv$.)
\end{proof}

El signo del producto escalar indica, por tanto, si dos vectores apuntan en direcciones parecidas u opuestas (\cref{tab:signo}). El ejercicio~\ref{eje:2-tres} explora hasta dónde pueden «llevarse mal» varios vectores del plano.

\begin{table}[h]
\centering\small
\caption{Signo del producto escalar y ángulo entre dos vectores no nulos.}\label{tab:signo}
\begin{tabular}{@{}lll@{}}
\toprule
Producto escalar & Ángulo & Interpretación \\
\midrule
$u\cdot v>0$ & $0\le\theta<\pi/2$ & direcciones parecidas \\
$u\cdot v=0$ & $\theta=\pi/2$ & direcciones perpendiculares \\
$u\cdot v<0$ & $\pi/2<\theta\le\pi$ & direcciones opuestas \\
\bottomrule
\end{tabular}
\end{table}

\begin{definicion}[Ortogonalidad]
Dos vectores $u,v\in\R^{d}$ son \emph{ortogonales}, y se escribe $u\perp v$, si $u\cdot v=0$.
\end{definicion}

El vector nulo es ortogonal a todos, y es el único vector ortogonal a sí mismo, por la positividad del producto escalar. Los vectores de la base canónica son ortogonales dos a dos: $e_i\cdot e_j=0$ si $i\neq j$.

\begin{teorema}[Pitágoras]\label{teo:pitagoras}
Si $u\perp v$, entonces $\norm{u+v}^{2}=\norm{u}^{2}+\norm{v}^{2}$.
\end{teorema}

\begin{proof}
$\norm{u+v}^{2}=\norm{u}^{2}+2\,u\cdot v+\norm{v}^{2}$, y el término central es nulo.
\end{proof}

\section{Proyección sobre una recta}

\begin{proposicion}[Proyección ortogonal sobre una recta]\label{prop:proyeccion}
Sean $u,a\in\R^{d}$ con $a\neq 0$. Entre todos los vectores de la recta $\{ta: t\in\R\}$, el más próximo a $u$ es
\[
p=\frac{u\cdot a}{a\cdot a}\,a,
\]
y es el único punto de la recta tal que $u-p\perp a$. El número $(u\cdot a)/\norm{a}$ se llama \emph{proyección escalar} de $u$ sobre $a$; si $\norm{a}=1$, coincide con $u\cdot a$.
\end{proposicion}

\begin{proof}
Sea $t^{*}=(u\cdot a)/(a\cdot a)$ y $p=t^{*}a$. Entonces $(u-p)\cdot a=u\cdot a-t^{*}\,a\cdot a=0$. Recíprocamente, $(u-ta)\cdot a=0$ obliga a $t=t^{*}$. Para cualquier $t$, el vector $u-ta=(u-p)+(t^{*}-t)a$ es suma de dos vectores ortogonales, y por el \cref{teo:pitagoras}
\[
\norm{u-ta}^{2}=\norm{u-p}^{2}+(t^{*}-t)^{2}\norm{a}^{2}\ge\norm{u-p}^{2},
\]
con igualdad solo si $t=t^{*}$.
\end{proof}

\begin{figure}[h]
\centering
\begin{tikzpicture}[scale=1.15,>={Stealth[length=2.4mm]},font=\small]
  \draw[thin,gray] (-0.4,-0.2) -- (5.2,2.6);
  \draw[->,thick] (0,0) -- (2.0,1.0) node[below right] {$a$};
  \draw[->,thick] (0,0) -- (2.6,2.8) node[above left] {$u$};
  \draw[->,thick,dashed] (0,0) -- (3.2,1.6) node[below right] {$p$};
  \draw[->] (3.2,1.6) -- (2.6,2.8) node[midway,right] {$u-p$};
  \draw (3.2,1.6) ++(-0.179,-0.089) -- ++(-0.089,0.179) -- ++(0.179,0.089);
  \node at (4.9,2.15) {$\{ta\}$};
\end{tikzpicture}
\caption{Proyección ortogonal de $u$ sobre la recta generada por $a$ (\cref{prop:proyeccion}). El error $u-p$ es perpendicular a $a$.}
\label{fig:proyeccion}
\end{figure}

\begin{ejemplo}\label{ej:proyeccion}
Sean $u=(4,-1,2)$ y $a=(1,1,1)$. Como $u\cdot a=5$ y $a\cdot a=3$, la proyección es $p=\tfrac53(1,1,1)$, y el error $u-p=\tfrac13(7,-8,1)$ cumple $(u-p)\cdot a=\tfrac13(7-8+1)=0$.
\end{ejemplo}

La proyección sobre una recta es el caso más sencillo de la proyección sobre un subespacio, que conduce a los mínimos cuadrados~\cite[§4.2]{strang}; en el \cref{tema:4} se escribe como producto de matrices (ejercicio~\ref{eje:4-proyeccion}).

\section{Similitud coseno}

\begin{definicion}[Similitud coseno]\label{def:coseno}
La \emph{similitud coseno} de dos vectores no nulos $u,v\in\R^{d}$ es
\[
\cos(u,v)=\frac{u\cdot v}{\norm{u}\,\norm{v}}=\hat u\cdot\hat v\in[-1,1].
\]
\end{definicion}

\begin{proposicion}\label{prop:coseno}
Sean $u,v\in\R^{d}$ no nulos.
\begin{enumerate}
  \item $\cos(au,bv)=\cos(u,v)$ para todos $a,b>0$, y $\cos(u,-v)=-\cos(u,v)$.
  \item $\cos(u,v)=1$ si y solo si $u=cv$ con $c>0$; $\cos(u,v)=-1$ si y solo si $u=cv$ con $c<0$.
\end{enumerate}
\end{proposicion}

\begin{proof}
(i) Numerador y denominador se multiplican por $ab=\abs{ab}$; con $b=-1$, el numerador cambia de signo y el denominador no. (ii) Si $\abs{\cos(u,v)}=1$, se da la igualdad en el \cref{teo:cs}, luego $u=cv$ con $c\neq 0$, y entonces $\cos(u,v)=c\norm{v}^{2}/(\abs{c}\norm{v}^{2})$ es el signo de $c$. El recíproco es inmediato.
\end{proof}

La similitud coseno mide, por tanto, solo la dirección: $\cos((1,2),(2,4))=1$ aunque los vectores tengan longitudes distintas.

\begin{ejemplo}[Embeddings de juguete]\label{ej:juguete}
Considérense vectores de $\R^{4}$ cuyas componentes se interpretan como (realeza, masculino, persona, comida):
\begin{center}\small
\begin{tabular}{@{}lcc@{\qquad}lcc@{}}
\toprule
Palabra & Vector & Norma & Palabra & Vector & Norma \\
\midrule
rey & $(2,1,2,0)$ & $3$ & hombre & $(0,1,2,0)$ & $\sqrt5$ \\
reina & $(2,-1,2,0)$ & $3$ & mujer & $(0,-1,2,0)$ & $\sqrt5$ \\
manzana & $(0,0,0,2)$ & $2$ & & & \\
\bottomrule
\end{tabular}
\end{center}
Se obtiene $\cos(\text{rey},\text{reina})=7/9\approx 0{,}7778$ (ángulo de $38{,}94^\circ$), $\cos(\text{rey},\text{hombre})=5/(3\sqrt5)\approx 0{,}7454$, $\cos(\text{rey},\text{mujer})=3/(3\sqrt5)\approx 0{,}4472$ y $\cos(\text{rey},\text{manzana})=0$. El vecino más próximo de «rey» es «reina»: comparten realeza y humanidad y solo discrepan en una componente. En un modelo real las componentes no tienen un significado asignado de antemano, pero la comparación se hace de la misma forma.
\end{ejemplo}

\begin{proposicion}[Coseno y distancia]\label{prop:coseno-distancia}
Si $\norm{u}=\norm{v}=1$, entonces $\norm{u-v}^{2}=2-2\cos(u,v)$. En consecuencia, entre vectores normalizados, ordenar los candidatos por distancia creciente a una consulta equivale a ordenarlos por similitud coseno decreciente.
\end{proposicion}

La demostración se propone en el ejercicio~\ref{eje:2-distancia}. Sin normalizar, ambas medidas pueden discrepar: un vector muy largo en la misma dirección que la consulta tiene coseno~1 pero puede estar lejos (ejercicio~\ref{eje:2-vecinos}). Los buscadores vectoriales de los sistemas RAG (módulo~22) usan el coseno o el producto escalar sobre vectores normalizados.

\section{Geometría en alta dimensión}

En $\R^{2}$, dos direcciones elegidas al azar forman ángulos de cualquier tamaño. En dimensión alta el comportamiento es muy distinto: dos direcciones aleatorias son casi ortogonales~\cite[§2.4.2]{bhk}. El resultado siguiente lo cuantifica.

\begin{proposicion}\label{prop:alta-dim}
Sea $U$ un vector aleatorio de $\R^{d}$ con $\norm{U}=1$ cuya distribución no cambia al permutar sus componentes ni al cambiar el signo de cualquiera de ellas (por ejemplo, $U=G/\norm{G}$ con $G_1,\dots,G_d$ normales estándar independientes). Para todo vector unitario fijo $v$,
\[
\E[U\cdot v]=0,\qquad \E\bigl[(U\cdot v)^{2}\bigr]=\frac1d .
\]
En particular, la similitud coseno entre $U$ y $v$ tiene media~0 y desviación típica $1/\sqrt d$.
\end{proposicion}

\begin{proof}
Cambiar el signo de la componente $i$ no altera la distribución, luego $\E[U_i]=\E[-U_i]$ y $\E[U_i]=0$. Del mismo modo, para $i\neq j$, $\E[U_iU_j]=\E[-U_iU_j]=0$. Por la invariancia bajo permutaciones, $\E[U_i^{2}]$ es el mismo número $c$ para todo $i$, y como $\sum_iU_i^{2}=1$, se tiene $dc=1$. Entonces $\E[U\cdot v]=\sum_iv_i\E[U_i]=0$ y
\[
\E\bigl[(U\cdot v)^{2}\bigr]=\sum_{i=1}^{d}\sum_{j=1}^{d}v_iv_j\,\E[U_iU_j]=\sum_{i=1}^{d}v_i^{2}\,\frac1d=\frac1d. \qedhere
\]
\end{proof}

Si $v$ también es aleatorio e independiente de $U$, el resultado se aplica condicionando al valor de $v$. Para $d=768$, la desviación típica del coseno entre dos direcciones aleatorias es $1/\sqrt{768}\approx 0{,}0361$; para $d=12\,288$, $\approx 0{,}0090$. Una simulación con 2000 pares de vectores gaussianos da $0{,}0363$ y $0{,}0091$, respectivamente. Por eso un coseno de $0{,}3$ entre dos embeddings de dimensión 768, que equivale a más de ocho desviaciones típicas, es una señal fuerte de parecido.

\begin{observacion}
Este fenómeno permite que un espacio de dimensión $d$ contenga muchas más de $d$ direcciones casi ortogonales entre sí. La hipótesis de \emph{superposición} de la interpretabilidad (módulo~23) se apoya en esta propiedad: un modelo puede representar más conceptos que dimensiones tiene su espacio, asignando a cada uno una dirección casi ortogonal a las demás.
\end{observacion}

\begin{lecturas}
\item Strang~\cite{strang}, §1.2, problemas sobre productos escalares y ángulos.
\item Deisenroth, Faisal y Ong~\cite{mml}, §3.3 «Lengths and Distances», §3.4 «Angles and Orthogonality» y §3.8 «Orthogonal Projections».
\item Boyd y Vandenberghe~\cite{vmls}, §3.2 «Distance» y §3.4 «Angle».
\item Blum, Hopcroft y Kannan~\cite{bhk}, §2.3 «The Geometry of High Dimensions» y §2.4.2 «Volume Near the Equator».
\end{lecturas}

\section*{Ejercicios del tema 2}
\addcontentsline{toc}{section}{Ejercicios del tema 2}

\begin{ejercicio}{1}{eje:2-escala}
Calcule el producto escalar y la similitud coseno entre $(1,2)$ y $(2,4)$, y entre $(1,2)$ y $(10,20)$. ¿Qué indica la comparación sobre lo que mide cada cantidad?
\end{ejercicio}

\begin{ejercicio}{1}{eje:2-vf}
Justifique si son verdaderas o falsas: (a) si $u\cdot v=0$, uno de los dos vectores es nulo; (b) $\cos(u,u)=1$ para todo $u\neq 0$; (c) si $\cos(u,v)=1$, entonces $u=v$.
\end{ejercicio}

\begin{ejercicio}{2}{eje:2-angulo}
Sean $u=(2,-1,2)$ y $v=(1,2,2)$. Calcule $u\cdot v$, $\norm{u}$, $\norm{v}$, $\cos(u,v)$ y el ángulo entre ambos en grados.
\end{ejercicio}

\begin{ejercicio}{2}{eje:2-vecinos}
Sean $a=(1,1)$, $b=(3,3)$ y $c=(1,-1)$. Calcule $\norm{a-b}$, $\norm{a-c}$, $\cos(a,b)$ y $\cos(a,c)$. ¿Qué vecino de $a$ elige cada medida? ¿Qué ocurre si antes se normalizan los tres vectores?
\end{ejercicio}

\begin{ejercicio}{2}{eje:2-normalizar}
Normalice $x=(3,4)$. Con $p=\hat x$ y $r=(1,0)$, compruebe numéricamente que $\norm{p-r}^{2}=2-2\cos(p,r)$.
\end{ejercicio}

\begin{ejercicio}{2}{eje:2-tres}
(a) ¿Pueden tres vectores $u,v,w\in\R^{2}$ cumplir a la vez $u\cdot v<0$, $v\cdot w<0$ y $u\cdot w<0$? Si es posible, dé un ejemplo con vectores unitarios y calcule el ángulo entre cada par. (b) ¿Pueden hacerlo cuatro vectores de $\R^{2}$, con los seis productos escalares negativos?
\fuente{adaptado de Strang~\cite{strang}, §1.2, problema~28, propuesto en el curso 18.06 del MIT~\cite{mit1806}, primavera de 2010, hoja de problemas~1.}
\end{ejercicio}

\begin{ejercicio}{3}{eje:2-distancia}
Demuestre la \cref{prop:coseno-distancia}: si $\norm{u}=\norm{v}=1$, entonces $\norm{u-v}^{2}=2-2\cos(u,v)$. Deduzca que, si $q,a,b$ son unitarios y $\cos(q,a)>\cos(q,b)$, entonces $\norm{q-a}<\norm{q-b}$. ¿Qué distancia corresponde a $\cos(u,v)=-1$?
\end{ejercicio}
